\documentclass[11pt,a4paper]{article}
% Shared preamble for RuPhO English problem statements (match T1 2026 style).
% Use: \input{latex/rupho_en_preamble.tex} after \documentclass.
\usepackage[margin=1in]{geometry}
\usepackage{amsmath,amssymb}
\usepackage{graphicx}
\usepackage{float}
\usepackage{enumitem}
\usepackage{microtype}
\usepackage{caption}
\usepackage{tabularx}
\usepackage{hyperref}

\captionsetup{font=small,labelfont=bf,skip=6pt}

\setlist[itemize]{leftmargin=1.2cm}
\setlist[enumerate]{leftmargin=1.2cm}

% One outer box per subproblem: [id | statement | points] (no inner rules)
\newcommand{\problembox}[3]{%
  \par\addvspace{0.42em}%
  \noindent
  \setlength{\fboxsep}{7.5pt}%
  \setlength{\fboxrule}{0.95pt}%
  \fbox{%
    \begin{minipage}{\dimexpr\linewidth-2\fboxsep-2\fboxrule\relax}
    \begin{tabularx}{\linewidth}{@{}>{\raggedright\arraybackslash\bfseries}p{2.1em} @{\hspace{0.9em}} >{\raggedright\arraybackslash}X @{\hspace{0.9em}} >{\raggedleft\arraybackslash\bfseries}p{2.4em}@{}}
    #1 & #2 & #3
    \end{tabularx}
    \end{minipage}%
  }%
  \par\addvspace{0.32em}%
}


\title{\textbf{Slipping rope}}
\author{\normalsize W24 (2024) --- English translation}
\date{}
\begin{document}
\maketitle

A smooth inclined plane makes an angle $\alpha$ with the horizontal. The surface of the inclined plane smoothly transforms into two weightless blocks $1$ and $2$ touching it, the radius of which can be considered negligibly small. A flexible rope with a constant linear mass density $\lambda$ is thrown over the blocks, the ends of which are initially located on the same horizontal line. The length of the part of the rope hanging on the left is $L_1$, and the length of the section of the inclined plane is $L$. Gravity acceleration $g$. Let us denote by $x$ the downward movement of the left end of the rope from the initial position (see figure). This problem is devoted to the study of the dynamics of a rope.

\begin{figure}[H]
    \centering
    \includegraphics[width=0.72\textwidth]{T1_EN_assets/fig1.jpg}
    \caption{Figure 1.}
\end{figure}

\section*{Part A. Weightless blocks (7.5 points)}

The rope is brought out of its equilibrium position by imparting low speed to it so that its left end begins to fall. Assume that the rope does not come off the blocks when moving.

\problembox{A1}{Determine the speed $v$ and acceleration $a$ of the left end of the rope at the moment when the left end is lowered by the amount $x$. Express your answer using $g$, $L_1$, $L$, $\alpha$ and $x$.}{1.50}

\problembox{A2}{Determine the amount of displacement of the left end of the rope $x_1$ at the moment in time when the tension of the section of the rope lying on the inclined plane is constant along the length of this section. Express your answer in terms of $L_1$, $L$ and $\alpha$.  At what ratios $L_1$, $L$ and $\alpha$ is the value $x_1$ less than the initial length of the right section of the rope?}{0.80}

\problembox{A3}{Determine the maximum tension force of the rope at the moment of time when its left end has shifted by a distance $x$. Express your answer in terms of $\lambda$, $g$, $L_1$, $L$, $\alpha$ and $x$.}{1.50}

The resulting expressions are applicable as long as the rope moves without being separated from the blocks. Now we will examine under what conditions the separation will occur.

\problembox{A4}{Determine the displacement values ​​$x_{{max}(1)}$ and $x_{{max}(2)}$ of the left end of the rope, at which it begins to lose contact with blocks $1$ and $2$, respectively, assuming that contact with the other block is maintained. Express your answers using $L_1$, $L$ and $\alpha$.}{2.50}

Next, assume that $L=2L_1$. The magnitude is $L=1~\text{м}$, and the acceleration due to gravity is $g=10~\text{м}/\text{с}^2$.

\problembox{A5}{Calculate the speed of the rope at the moment of its separation from one of the blocks if the inclined plane forms angles $\alpha_1=30^{\circ}$, $\alpha_2=45^{\circ}$ and $\alpha_3=60^{\circ}$ with the horizon.}{1.20}

\section*{Part B. Motion with a spring and a load (4.5 points)}

In this part of the problem, a mass of mass $m$ is attached to the left end of the rope, and a weightless spring is attached to the right end of the rope. The spring under consideration begins to stretch only if a force greater than a certain minimum value $F_0$ is applied to its ends, and it cannot compress. A qualitative graph of the dependence of the elastic force of a spring on its elongation is shown in the figure below. The value $F_0$ ($F_0 < mg$) and the spring stiffness coefficient $k$ ($k = \Delta F/\Delta l$) are considered given. The second end of the spring is fixed, the spring is oriented vertically. The geometric dimensions of the rope and its initial position are the same as in part A.

\begin{figure}[H]
    \centering
    \includegraphics[width=0.72\textwidth]{T1_EN_assets/fig2.jpg}
    \caption{Figure 2.}
\end{figure}

\begin{figure}[H]
    \centering
    \includegraphics[width=0.72\textwidth]{T1_EN_assets/fig3.jpg}
    \caption{Figure 3.}
\end{figure}

\problembox{B1}{Determine the displacement value $x_0$ of the left end of the rope corresponding to the equilibrium position of the system. Express your answer in terms of $m$, $\lambda$, $L_1$, $L$, $\alpha$, $g$, $k$ and $F_0$.}{0.80}

\problembox{B2}{The equilibrium position from the previous paragraph exists if the spring stiffness exceeds some minimum value $k_{min}$. Identify $k_{min}$. Express your answer using $m$, $\lambda$, $L_1$, $L$, $\alpha$, $g$, $k$ and $F_0$. Further, assume everywhere that $k > k_{min}$.}{0.70}

Let us denote the potential energy of the load and rope in the gravitational field as $W_{p(\text{г})}$ and $W_{p(\text{в})}$, respectively, and the potential energy of deformation of the spring as $W_{p(\text{пр})}$. The potential energy $W_p$ of the entire system is their sum: $$W_p=W_\{p(\textbackslash\{\}text\{г\})\}+W_\{p(\textbackslash\{\}text\{в\})\}+W_\{p(\textbackslash\{\}text\{пр\})\} $$

\problembox{B3}{Find the change in the potential energy of the system $\Delta{W}_p$ when the left end of the rope is displaced by the amount $x$. Express your answer in terms of $m$, $\lambda$, $L_1$, $L$, $\alpha$, $g$, $k$, $F_0$ and $x$.}{0.80}

\problembox{B4}{If we take the potential energy $W_p$ at the equilibrium position equal to zero, then it can be represented in the form $$W_p=\textbackslash\{\}cfrac\{A(x-B)^2\}\{2\}. $$Определите $A$ и $B$. Ответы выразите через $m$, $\lambda$, $L_1$, $L$, $\alpha$, $g$, $k$ и $F_0$.}{0.70}

The parameters of the system are such that if it is released from the initial position without an initial speed, then the right end of the thread does not reach the level of block $2$.

\problembox{B5}{Determine the maximum speed of the load $v_{max}$ and its displacement from the initial position $x_{max}$ if the system is released without an initial speed.  Express your answers using $m$, $\lambda$, $L_1$, $L$, $\alpha$, $g$, $k$ and $F_0$.}{0.50}

\problembox{B6}{What minimum speed $v_{min}$ must the weight have in the initial position for the right end of the thread to reach the level of block $2$? Express your answer using $m$, $\lambda$, $L_1$, $L$, $\alpha$, $g$, $k$ and $F_0$.}{1.00}

\vspace{1em}
\noindent\textit{Source:} \url{https://pho.rs/p/3575}

\end{document}